% occurrence_lemma.tex -- the load-bearing lemma for the L+lim universality
% construction (R3, item 1).  Self-contained; written in the notation of
% docs/proof/main.tex ([A/B] substitution, \sigma/\tau alphabets).
% Machine verification: research/scratch/lim/verify_r3.py, part A.

\subsection{The occurrence lemma for the two-counter compilation}

Fix $\Sigma=\{\sigma,\tau\}$ (any two characters of the ambient
alphabet; the compiler is symmetric under swapping them).  Write
$V(k)$ for the \emph{frame family}
\[
V(k) \;=\; \sigma^2(\tau\sigma)^{k-1}\tau\sigma^3 ,
\qquad k \ge 1 ,
\]
i.e.\ the binary string $\sigma\sigma\,\tau\sigma\,\tau\sigma\cdots
\tau\sigma\,\sigma\sigma$ with $k$ copies of $\tau\sigma$, of length
$2k+4$.  $V(k)$ has $k$ isolated $\tau$'s, head run $\sigma^2$, $k-1$
interior $\sigma$-runs of length $1$, and tail run $\sigma^3$;
distinct $k$ give pairwise non-occurring strings (the $\sigma$-run
signature $[2,1^{\,k-1},3]$ forces any occurrence to consume the
whole string, as proved below).

\paragraph{Configurations.}
A compiled two-counter machine with instructions $1,\dots,s$ and halt
slot $H$ uses \emph{frame constants}: delimiters $D_0,D_1,D_2,D_3$,
instruction codes $P_1,\dots,P_s,P_H$, and the wrap marker $M$ -- each
assigned a distinct $V(k)$ with $k\ge 2$.  A configuration is
\[
\cfg \;=\; D_0\;\Pi\;D_1\;\tau^{x}\;D_2\;\tau^{y}\;D_3 ,
\]
where the \emph{program zone} $\Pi$ lists $P_1\cdots P_s\,P_H$ in
order with the \emph{current} instruction wrapped between two copies
of $M$: $\cdots M P_i M\cdots$ (write $\pc=i$; write $\pc=0$ for the
halted configuration $P_1\cdots P_s\,M P_H M$), and where
$\tau^{x},\tau^{y}\ge 0$ are the \emph{tallies} of the two counters
(an empty tally makes its delimiters adjacent).  The step function
and the output stage are built from the following \emph{patterns}:
the dispatch/unwrap patterns $M P_i M$, the wrap targets $P_j$, the
increment patterns $D_r$, the decrement patterns $D_r\tau$, the zero
tests $D_rD_{r'}$ (adjacent delimiters), the halt dispatch $M P_H M$,
and the output prefix $D_0\Pi(0)D_1D_2$.

\begin{lemma}[occurrence]\label{lem:occurrence}
Let the frame constants be assigned pairwise distinct family indices
$k\ge 2$.  Then for every $\pc\in\{0,\dots,s\}$ and every $x,y\ge 0$,
in $\cfg(\pc,x,y)$:
\begin{enumerate}
\item[(i)] every frame constant occurs exactly at its slots ($M$
      twice, at the two wrap positions; every other constant once);
\item[(ii)] $M P_i M$ occurs iff $\pc=i$ ($i\in\{1,\dots,s\}$), and
      $M P_H M$ occurs iff $\pc=0$; in either case exactly once;
\item[(iii)] $D_rD_{r'}$ occurs iff counter $r$ is zero, exactly once;
\item[(iv)] $D_r\tau$ occurs iff counter $r$ is nonzero, exactly
      once, and its occurrence ends at the \emph{first} tally
      character;
\item[(v)] the output prefix $D_0\Pi(0)D_1D_2$ occurs iff $\pc=0$ and
      $x=0$, exactly once, at position $0$.
\end{enumerate}
\end{lemma}

\begin{proof}
The proof is a \emph{gap calculus}.  For a string $W$ containing at
least one $\tau$ and no $\tau\tau$, let its \emph{gap sequence} be
$(h;\,g_1,\dots,g_{m-1};\,t)$ where $h$ is the number of $\sigma$'s
before the first $\tau$, $g_i$ the number of $\sigma$'s between the
$i$-th and $(i{+}1)$-st $\tau$, and $t$ the number of $\sigma$'s after
the last $\tau$.

\emph{Step 1 (the gap sequences of configurations).}  Every frame
constant has gap sequence $(2;\,1^{\,k-1};\,3)$.  Concatenating two
frame constants merges the trailing $\sigma^3$ of the first with the
leading $\sigma^2$ of the second: an interior gap of $5$.  A tally
$\tau^{n}$ ($n\ge 1$) between its delimiters contributes the gap
pattern $3\,0^{\,n-1}\,2$ (delimiter tail $\sigma^3$, then $n-1$ gaps
of $0$ inside the tally, then delimiter head $\sigma^2$); an empty
tally makes the delimiters touch, a gap of $5$.  So the gap sequence
of any configuration is: head $2$ (the string starts with $D_0$'s
$\sigma^2$), then the constants' interior $1$-blocks interleaved with
junction gaps of $5$, with the two tally junctions $3\,0^{x-1}\,2$
(or $5$), and tail $3$ (the string ends with $D_3$'s $\sigma^3$).

\emph{Step 2 (patterns cannot cross tallies).}  Every pattern has at
least two $\tau$'s, no $\tau\tau$, head exactly $2$, tail in
$\{0,3\}$, and interior gaps in $\{1,3,5\}$ only -- never $0$ and
never $2$.  Now let a pattern $W$ occur in $\cfg$.  Matching
character-for-character, each $\tau$ of $W$ is a $\tau$ of $\cfg$,
and since between two consecutive $\tau$'s of $W$ there are only
$\sigma$'s, the $\tau$'s of $W$ map to \emph{consecutive} $\tau$'s of
$\cfg$ (a skipped $\tau$ would have to appear among $W$'s
$\sigma$'s).  Hence every interior gap of $W$ equals the
corresponding gap of $\cfg$.  Gaps of $0$ occur in $\cfg$ only
inside tallies of length $\ge 2$, so $W$ covers no two consecutive
tally characters.  A gap of $2$ occurs in $\cfg$ only after the last
character of a tally, so $W$'s $\tau$-sequence cannot cross a tally's
right edge; and if $W$'s \emph{first} $\tau$ were a tally character,
then $W$'s next gap would have to be $0$ (deeper into the tally) or
$2$ (out of it) -- both impossible -- unless the tally has length
exactly $1$ and $W$'s next gap is $2$: also impossible.  The only
pattern whose $\tau$-sequence can meet a tally at all is one whose
\emph{last} gap is $3$ -- the decrement pattern $D_r\tau$ -- and the
gap $3$ pins that last $\tau$ to the tally's \emph{first} character
(the gap $3$ occurs in $\cfg$ only at a nonzero tally's left edge).
This proves (iv) modulo the pinning of Step 3, and shows that every
other pattern avoids tallies entirely.

\emph{Step 3 (pinning).}  It remains to show that a pattern that
avoids tallies aligns slot-exactly with the frame constants.  Let
$W$'s leading segment (up to its first gap of $3$ or $5$) contain
$k$ $\tau$'s, all separated by gaps of $1$.  The first of these
$\tau$'s must be a constant's \emph{first} $\tau$: a middle $\tau$
of a constant has only one $\sigma$ before it, but $W$'s head is
$\sigma^2$ and the head run must fit inside the preceding
$\sigma$-run.  If the constant had more than $k$ $\tau$'s, the
$(k{+}1)$-st would follow a gap of $1$, where $W$ has its seam $3$
or $5$ -- impossible; so the constant has exactly $k$ $\tau$'s, and
by the distinctness of family indices it is the unique slot holding
that constant string.  The same argument applies constant-by-constant
along $W$'s seams (each seam of $5$ matches a junction, i.e.\ the
$\tau$'s on both sides belong to adjacent slots; the seam of $3$
matches only a tally's left edge).  Assembling: $M P_i M$ can only
chain $M$-slot, junction, $P_i$-slot, junction, $M$-slot, and the two
$M$ slots in a configuration bracket exactly one code -- the wrapped
one -- giving (ii).  $D_rD_{r'}$ chains $D_r$'s slot to the
\emph{next} slot only through a gap of $5$, which after $D_r$'s slot
arises exactly when the tally is empty: (iii).  The output prefix is
pinned at $D_0$'s slot and then slot-by-slot through the program
zone (which matches $\Pi(0)$ iff $\pc=0$), then across the first
tally junction ($5$ iff $x=0$): (v).  Finally (i): a single frame
constant $V(k)$ is itself a pattern with leading $1$-block of length
$k-1\ge 1$, pinned by Step 3 to a slot holding that string; the
slots holding $V(k_M)$ are exactly the two wrap positions of $M$,
every other constant has exactly one slot.
\end{proof}

\begin{remark}
The hypothesis $k\ge 2$ is used only to keep every pattern at two or
more $\tau$'s (a lone $V(1)$ would be pinned only by its neighbors).
The machine verification (part A of \texttt{verify\_r3.py}) also runs
the full census under the hypothesis-violating assignment $k=1$ for
one constant and finds no spurious occurrence there either -- the
\emph{sequence} of constants in a long pattern pins it even when an
individual constant is not unique -- but the clean statement above
does not need this.
\end{remark}

\begin{remark}[why the census cap cannot hide a failure]
Lemma~\ref{lem:occurrence} is a statement about \emph{all}
$x,y\ge 0$, and its proof is the finite case analysis of Steps 1--3:
the only $x,y$-dependent gaps are the tally patterns
$3\,0^{\,n-1}\,2$, and the case split on $n=0$, $n=1$, $n\ge 2$ is
exhaustive.  The machine check therefore samples edge lengths
$x,y\in\{0,1,2\}$ at every $\pc$ explicitly, all $x,y\le 40$
systematically, and $400$ random states up to $1000$: the census
counts every pattern in every configuration and agrees with the
lemma's predictions everywhere.
\end{remark}
